% !TeX encoding = UTF-8
\documentclass[variant=docente,cover=institutional,mode=screen]{ua-engineering-v6-article}
% !TeX encoding = UTF-8
\UASetUniversity{Universidad Austral}
\UASetFaculty{Facultad de Ingeniería}
\UASetDepartment{Departamento de Computación}
\UASetCampus{Pilar}
\UASetCareer{Ingeniería Informática}
\UASetCommission{Comisión A}
\UASetProfessor{Equipo docente}
\UASetSemester{Segundo cuatrimestre 2026}
\UASetCourse{Sistemas Distribuidos}
\UASetDate{2026}


\UASetClassNumber{01}
\UASetClassTitle{Pensamiento distribuido}
\UASetDocKind{Guía docente}

\title{Clase 01 --- Pensamiento distribuido}
\author{\UAProfessor}
\date{\UADate}

\begin{document}
\maketitle

\section{Objetivo pedagógico profundo}\label{sec:objetivo-pedagogico-profundo}

La meta de esta clase es producir una ruptura conceptual.
El estudiante debe dejar de razonar como si el sistema tuviera un estado global único, inmediatamente observable y perfectamente interpretable.

La clase no debe medirse por cuántos términos nuevos aparecen, sino por si el estudiante logra aceptar esta afirmación:

\begin{center}
\textbf{“Desde un borde remoto del sistema, a veces no puedo saber qué pasó realmente.”}
\end{center}

Si esa idea no queda instalada, el resto del curso se vuelve una colección de técnicas sin fundamento.

\section{Resultado de aprendizaje esperado}\label{sec:resultado-de-aprendizaje-esperado}

Al terminar la clase, el estudiante debería poder:

\begin{itemize}
\item Explicar qué vuelve distribuido a un problema;
\item construir dos historias indistinguibles para un observador remoto;
\item distinguir síntoma observado de causa subyacente;
\item documentar observaciones, hipótesis, evidencias y conclusiones;
\item relevar, analizar y proponer separando cada dimensión y conectándolas en una decisión coherente;
\item usar el esquema $D=(P,F,G,S,T)$ para ordenar una discusión técnica;
\item leer CAP como una restricción bajo partición.
\end{itemize}

\section{Duración sugerida}\label{sec:duracion-sugerida}

Clase de 90 minutos.

\section{Estructura sugerida}\label{sec:estructura-sugerida}

\subsection*{Bloque 0 (10 min): presentación de la materia}

Presentar brevemente el recorrido de la asignatura, el plantel docente y el esquema de evaluación.

Usar este bloque para instalar desde el inicio las habilidades esperadas:

\begin{itemize}
\item relevar contextos de sistemas distribuidos;
\item analizar síntomas, restricciones y garantías;
\item proponer y defender alternativas de diseño con trade-offs explícitos.
\end{itemize}

La presentación administrativa no debe quedar separada del objetivo pedagógico: sirve para mostrar que la materia evalúa razonamiento de ingeniería, no memorización de tecnologías.

\subsection*{Bloque 1 (10--15 min): abrir con un caso, no con una definición}

No empezar diciendo “un sistema distribuido es\ldots”.  
Empezar con el ejemplo del pago que no devuelve respuesta.

Preguntar al curso:

\begin{quote}
¿El pago se realizó o no?
\end{quote}

Lo esperable es que haya respuestas rápidas e intuitivas.
Esa intuición debe ser tensionada.

\subsection*{Bloque 2 (15--20 min): construir la indistinguibilidad}

Presentar dos historias:

\begin{itemize}
\item El servidor nunca recibió la solicitud;
\item el servidor sí procesó el pago, pero la respuesta se perdió.
\end{itemize}

Preguntar:

\begin{quote}
¿Qué ve el cliente en cada caso?
\end{quote}

El objetivo es que el grupo vea que la observación puede ser la misma aunque la historia interna sea distinta.

\subsection*{Bloque 3 (15 min): formalizar qué vuelve distribuido a un problema}

Recién después del shock conceptual, introducir la definición operativa.

Explicar que lo distribuido no es solo “varias máquinas”, sino coordinación bajo observación parcial, sin reloj global y con comunicación no confiable.

\subsection*{Bloque 4 (15 min): sincronía y asincronía}

Mostrar que la intuición de muchos estudiantes está más cerca de un modelo síncrono ideal que de un sistema realista.

El objetivo no es hacer teoría de modelos en profundidad todavía, sino dejar claro por qué una demora arbitrariamente larga ya destruye varias intuiciones locales.

\subsection*{Bloque 5 (15 min): framework del curso}

Presentar:

\[
D = (P, F, G, S, T)
\]

y usarlo inmediatamente sobre el caso del pago.

Este momento es clave porque instala una disciplina intelectual: no empezar por la solución, sino por el problema, las fallas y la garantía.

\subsection*{Bloque 6 (10 min): cierre con CAP}

Usar CAP para cerrar la clase de forma fuerte, pero correcta.

No enseñar CAP como trivia.
Primero presentar el teorema en pocas palabras: en un sistema replicado, bajo partición, no se pueden garantizar simultáneamente consistencia fuerte, disponibilidad total y tolerancia a particiones.
Después remarcar la interpretación madura: como las particiones son una falla que el sistema debe tolerar, la decisión real es entre preservar consistencia fuerte o preservar disponibilidad total mientras dura la partición.

\section{Qué explicar y por qué}\label{sec:que-explicar-y-por-que}

\subsection*{1. El ejemplo del pago}
Debe explicarse por qué es accesible, realista y rompe rápido la ilusión de certeza.

\subsection*{2. Indistinguibilidad}
Debe explicarse por qué es la idea fundacional del curso.
Sin esto, el alumno no entiende por qué aparecen retries, idempotencia, consenso o CAP\@.

\subsection*{3. Modelo asíncrono}
Debe explicarse por qué fuerza a abandonar la intuición de que el tiempo siempre sirve para distinguir causas.

\subsection*{4. Framework del curso}
Debe explicarse por qué ordena toda la materia y evita que la discusión tecnológica tape el problema conceptual.

\subsection*{5. CAP}
Debe explicarse por qué da cierre y muestra que las restricciones del curso no son caprichosas, sino estructurales.

\section{Qué no hacer}\label{sec:que-no-hacer}

\begin{itemize}
\item No abrir la clase con una definición enciclopédica.
\item No usar CAP como slogan.
\item No apresurarse a introducir algoritmos o tecnologías.
\item No resolver la incertidumbre demasiado rápido con “ponemos un timeout y listo”.
\end{itemize}

\section{Preguntas disparadoras recomendadas}\label{sec:preguntas-disparadoras-recomendadas}

\begin{itemize}
\item ¿Qué sabe realmente el cliente?
\item ¿Qué parte de tu respuesta depende de una suposición no garantizada?
\item ¿Qué historia interna distinta podría ser compatible con la misma observación?
\item ¿Qué garantía querés preservar en este caso concreto?
\item ¿Qué estarías dispuesto a sacrificar para preservarla?
\end{itemize}

\section{Errores frecuentes del alumnado}\label{sec:errores-frecuentes-del-alumnado}

\begin{itemize}
\item Creer que timeout implica crash;
\item pensar que siempre puede saberse “qué pasó realmente”;
\item asumir que toda garantía importante puede obtenerse sin costo;
\item usar CAP como frase de moda;
\item creer que el curso trata sobre herramientas y no sobre decisiones bajo incertidumbre.
\end{itemize}

\section{Momento crítico de la clase}\label{sec:momento-critico-de-la-clase}

El momento crucial ocurre cuando el estudiante deja de insistir en responder “sí” o “no” a la pregunta del pago y acepta la respuesta correcta:

\begin{center}
\textbf{“No puedo saberlo con certeza desde acá.”}
\end{center}

Ese momento debe señalarse explícitamente.
Si el alumno no atraviesa esa ruptura, la clase puede parecer interesante, pero no transforma su forma de razonar.

\section{Criterio de éxito docente}\label{sec:criterio-de-exito-docente}

La clase salió bien si el estudiante puede:

\begin{itemize}
\item Construir dos historias indistinguibles;
\item explicar por qué una observación no identifica por sí sola una causa;
\item usar el esquema $D=(P,F,G,S,T)$ sin mezclar niveles;
\item defender una garantía junto con el trade-off que requiere.
\end{itemize}

\begin{uawarn}
No evaluar esta clase por memorización de términos, sino por cambio de modelo mental.
\end{uawarn}

\section{Conexión con la clase 2}\label{sec:conexion-con-la-clase-2}

La clase 1 debe dejar preparada la siguiente pregunta:

\begin{quote}
Si no puedo saber exactamente qué pasó, ¿cómo diseño comunicación remota razonable?
\end{quote}

Esa pregunta abre naturalmente la clase 2 sobre:
\begin{itemize}
\item RPC,
\item timeout,
\item retry,
\item idempotencia.
\end{itemize}


\section{Plan de conducción ampliado}
La clase debe alternar explicación, predicción y contraste. Antes de revelar cada mecanismo, pedir al grupo que anticipe qué puede salir mal; después, volver al mismo escenario y preguntar qué propiedad mejora y qué costo aparece. Cada bloque conceptual debe cerrar con una decisión de diseño observable.
\subsection*{Secuencia recomendada}
\begin{enumerate}
\item Recuperar el prerequisito de la clase anterior.
\item Presentar un caso mínimo que falle y solicitar hipótesis.
\item Formalizar el concepto central de Pensamiento distribuido, distinguiendo seguridad, progreso y supuestos.
\item Resolver un ejemplo completo en pizarra, incluyendo una falla intermedia.
\item Comparar una alternativa de diseño y explicitar trade-offs.
\item Cerrar con una pregunta del Trabajo Final y una pregunta de defensa oral.
\end{enumerate}
\section{Diagnóstico de comprensión durante la clase}
No preguntar solamente ``¿se entendió?''. Usar preguntas discriminantes: ``¿qué cambia si el mensaje llega dos veces?'', ``¿qué propiedad sigue siendo cierta si cae un nodo?'', ``¿esto demuestra seguridad o progreso?'', ``¿qué observa un cliente externo?''. Si el estudiante responde solo con nombres de patrones, pedir una ejecución concreta.
\section{Criterios para corregir ejercicios}
Una respuesta excelente declara supuestos, construye una secuencia verificable, nombra la garantía con precisión, reconoce un costo y conecta la decisión con el dominio. Penalizar afirmaciones absolutas sin condiciones.
\section{Vinculación con el Trabajo Final Integrador}
Reservar entre 5 y 10 minutos para que cada grupo registre una decisión con: problema, alternativa elegida, alternativa descartada y razón. Esto crea trazabilidad entre las clases y las entregas acumulativas.
\section{Lectura docente previa}
Revisar la bibliografía indicada en los apuntes y preparar al menos un contraejemplo que no aparezca literalmente en las diapositivas. El objetivo es poder responder preguntas de frontera sin convertir la clase en una lectura del deck.

\section{Arquitectura didáctica v9 para 180 minutos}
\subsection*{Bloque 1 - desafío inicial (20 min)}
Presentar una situación donde aparezca delay, pérdida, duplicación, partición y observación incompleta. Pedir predicciones individuales antes de introducir vocabulario. El objetivo es hacer visibles los supuestos intuitivos y recuperar el prerequisito de la clase anterior.

\subsection*{Bloque 2 - construcción conceptual (35 min)}
Formalizar incertidumbre, fallas parciales, asincronía, indistinguibilidad y marco P-F-G-S-T. Separar seguridad, progreso y experiencia del usuario. Explicar no solo \emph{qué} hace cada mecanismo, sino \emph{por qué} existe y qué supuesto permite que funcione.

\subsection*{Bloque 3 - modelo y figura (35 min)}
Construir en pizarra una línea temporal, grafo, log, máquina de estados o tabla de decisiones. Recorrer una ejecución nominal y una adversarial. La representación debe quedar como referencia para los apuntes y el ejercicio principal.

\subsection*{Pausa (10 min)}
Recoger una pregunta anónima o un punto de confusión. Elegir una pregunta que permita distinguir síntoma, causa y propiedad.

\subsection*{Bloque 4 - caso trabajado con falla (40 min)}
Aplicar el mecanismo al caso conductor. Introducir una falla a mitad de ejecución y detenerse en cada transición: quién sabe qué, qué se persiste, qué garantía sigue vigente y qué observa el cliente.

\subsection*{Bloque 5 - taller de decisión (25 min)}
Los grupos aplican P-F-G-S-T a su Trabajo Final. Deben producir un ADR breve con alternativa elegida, alternativa descartada y evidencia pendiente. La evidencia de esta clase es timeline de ejecuciones indistinguibles y explicación de qué sabe cada actor.

\subsection*{Bloque 6 - cierre y exit ticket (15 min)}
Cada estudiante responde: propiedad principal, supuesto decisivo, costo aceptado y condición bajo la cual cambiaría de diseño. Usar las respuestas para iniciar la clase siguiente.

\section{Qué explicar, por qué y cómo verificarlo}
\begin{itemize}
\item Explicar el problema antes del producto, porque reconocer la familia de problema es el resultado cognitivo central.
\item Explicar el invariante, porque permite evaluar configuraciones y no solo repetir un flujo nominal.
\item Explicar el límite, porque delay, pérdida, duplicación, partición y observación incompleta puede invalidar supuestos o reducir progreso.
\item Verificar comprensión mediante timeline de ejecuciones indistinguibles y explicación de qué sabe cada actor, no mediante una definición aislada.
\end{itemize}

\section{Preguntas socráticas}
\begin{itemize}
\item ¿Qué sabe cada actor en este punto de la ejecución?
\item ¿Qué propiedad sigue siendo cierta si la respuesta nunca llega?
\item ¿Qué parte es safety y cuál es liveness?
\item ¿Qué alternativa más simple resolvería el problema si relajamos una garantía?
\item ¿Qué observación refutaría nuestra hipótesis de diseño?
\end{itemize}

\section{Errores previsibles y estrategia de corrección}
\begin{itemize}
\item Si el alumno responde con un nombre de tecnología, pedir actores, mensajes y estados.
\item Si confunde timeout con falla, construir dos ejecuciones indistinguibles.
\item Si promete todas las garantías, pedir comportamiento bajo partición o saturación.
\item Si mide solo rendimiento, preguntar cómo evidencia la propiedad semántica.
\item Si la solución es excesiva, pedir la alternativa mínima y el costo marginal de la garantía fuerte.
\end{itemize}

\section{Criterios de evaluación formativa}
Una respuesta excelente declara supuestos, recorre una ejecución, nombra con precisión la garantía, reconoce el costo y propone evidencia reproducible. Una respuesta suficiente identifica el mecanismo y el trade-off principal. Una respuesta insuficiente usa slogans, omite fallas o confunde producto con propiedad.

\section{Preparación docente y bibliografía}
Lectura previa recomendada: Kleppmann, introducción y capítulo 8; Tanenbaum y van Steen, modelos de sistema. Preparar un contraejemplo adicional y una variante del caso en la que cambie el costo del error; esto evita convertir la clase en una lectura de diapositivas.

\end{document}
